\documentclass[11pt]{article}

\usepackage[margin=1in]{geometry}
\usepackage{amsmath,amssymb,amsthm,mathtools}
\usepackage{bm}
\usepackage{booktabs}
\usepackage{array}
\usepackage{longtable}
\usepackage{hyperref}
\usepackage{microtype}
\usepackage[T1]{fontenc}
\usepackage[utf8]{inputenc}

\title{A Cohomological Statistical Mechanics of Go:\\
Local Constraints, Rooted Forests, Loop Gases, and Emergent Enclosure}
\author{Zhengbang Zhou}
\date{}

\newtheorem{proposition}{Proposition}
\newtheorem{conjecture}{Conjecture}

\begin{document}
\maketitle

\begin{abstract}
The game of Go suggests an unusual statistical-mechanics problem. Its microscopic degrees of freedom are entirely local---a site is Black, White, or empty, and stones interact through nearest-neighbor connectivity---yet its meaningful macroscopic objects are intrinsically nonlocal: groups, liberties, eyes, enclosure, territory, and life. The central question of this note is whether these global structures can be understood as emergent geometric and topological observables of a locally constrained lattice system.

A particularly useful formulation starts from a spin-1 lattice gas,
\[
s_i\in\{-1,0,+1\},
\]
and augments an ordinary Blume--Emery--Griffiths-type Hamiltonian with an exact representation of the Go liberty constraint. For each color, the graph of same-color stones has a graph Laplacian \(L_c\). Stones adjacent to empty sites receive a diagonal root weight \(R_c\). Then
\[
\det(L_c+R_c)>0
\]
if and only if every connected group of color \(c\) possesses at least one liberty. Moreover, this determinant is the partition function of rooted spanning forests. Thus the rule ``every group must reach a liberty,'' which appears globally nonlocal, admits a local auxiliary representation in terms of directed forests and, after integrating those degrees of freedom out, becomes a Laplacian determinant.

The resulting model sits at an intersection of several established areas of mathematical physics: spin-1 lattice gases, Peierls contour gases, \(O(n)\) loop models, bootstrap and kinetically constrained dynamics, random-cluster models, spanning forests, fermionic/supersymmetric field theories, algebraic topology, and persistent homology. The interesting claim is therefore not that Go is an Ising model. The more ambitious proposal is that Go provides a natural example of a locally constrained system in which complement topology---inside versus outside, enclosed chambers, and eventually territory---becomes a dynamically meaningful macroscopic degree of freedom.
\end{abstract}

\section{What exactly are we trying to explain?}

There are three progressively stronger questions.

The weakest question is whether a lattice Hamiltonian can assign plausible energies to Go positions. That has essentially been done in various Ising-inspired and graphical-model approaches.

The more interesting question is whether one can encode the actual structural constraint of Go groups rather than merely producing a heuristic evaluation function.

And the strongest question is
\[
\boxed{
\text{Can local microscopic constraints generate a macroscopic notion of enclosure?}
}
\]

The desired chain is
\[
\text{local adjacency}
\longrightarrow
\text{connectivity}
\longrightarrow
\text{survival constraint}
\longrightarrow
\text{extended barriers}
\longrightarrow
\text{components of the complement}
\longrightarrow
\text{enclosed area}.
\]

The central physics question then becomes whether there is a regime in which these enclosed regions become macroscopic and whether the onset of such enclosure behaves like a genuine phase transition.

\section{The microscopic state space}

Let
\[
\Lambda_L=\{1,\ldots,L\}^2
\]
be an \(L\times L\) square lattice with nearest-neighbor edges. The physical Go board corresponds to \(L=19\), but statistical mechanics requires allowing \(L\to\infty\).

Associate to every site \(i\in\Lambda_L\) a spin-1 variable
\[
s_i\in\{-1,0,+1\},
\]
with
\[
s_i=
\begin{cases}
+1,&\text{Black stone},\\
-1,&\text{White stone},\\
0,&\text{empty}.
\end{cases}
\]

Two useful derived variables are
\[
n_i=s_i^2
\]
and
\[
m_i=s_i.
\]

The occupation variable \(n_i\in\{0,1\}\) says whether the site contains a stone, while \(m_i\) carries the color.

The corresponding projectors are
\[
P_i^B=\frac{s_i^2+s_i}{2},
\qquad
P_i^W=\frac{s_i^2-s_i}{2},
\qquad
P_i^0=1-s_i^2.
\]

\section{The natural unconstrained Hamiltonian: Blume--Emery--Griffiths}

The most natural standard condensed-matter starting point is not the ordinary Ising model but the spin-1 Blume--Emery--Griffiths family,
\[
\boxed{
H_0[s]
=
-J\sum_{\langle ij\rangle}s_i s_j
-K\sum_{\langle ij\rangle}s_i^2s_j^2
+\Delta\sum_i s_i^2
-h\sum_i s_i.
}
\]

For two adjacent stones of the same color,
\[
s_i s_j=+1,
\]
so their pair energy is
\[
-J-K.
\]

For neighboring stones of opposite colors,
\[
s_i s_j=-1,
\]
giving
\[
+J-K.
\]

A stone-empty pair receives neither interaction.

Thus \(J\) controls color alignment, while \(K\) controls occupation clustering independently of color. The crystal-field parameter \(\Delta\) is a chemical potential for putting stones on the board, and \(h\) produces a Black/White bias.

\section{Groups and liberties}

For a fixed color \(c\in\{+1,-1\}\), define
\[
V_c(s)=\{i:s_i=c\}.
\]

Let \(G_c(s)\) be the graph whose vertices are \(V_c(s)\) and whose edges join nearest-neighbor stones of the same color. Connected components of \(G_c\) are exactly Go groups.

For a stone \(i\), define its number of immediately adjacent empty sites,
\[
\ell_i(s)
=
\sum_{j\sim i}P_j^0.
\]

A group \(C\subseteq V_c\) has at least one liberty exactly when
\[
\sum_{i\in C}\ell_i>0.
\]

The important point is that the local quantity \(\ell_i\) is easy to compute, but
\[
\boxed{
\text{Every connected component contains some }i\text{ with }\ell_i>0
}
\]
is a global graph property.

\section{Turning the liberty rule into a Laplacian determinant}

Let \(L_c\) be the weighted graph Laplacian of \(G_c\). For simplicity take unit edge weights:
\[
(L_c)_{ij}
=
\begin{cases}
\deg_c(i),&i=j,\\
-1,&i\sim j,\ s_i=s_j=c,\\
0,&\text{otherwise}.
\end{cases}
\]

Now define a diagonal matrix
\[
R_c=\operatorname{diag}(r_i),
\]
where
\[
r_i=\lambda\,\ell_i,
\qquad \lambda>0.
\]

A stone adjacent to one or more liberties therefore receives positive root weight; a stone with no adjacent liberty receives \(r_i=0\).

Define
\[
\boxed{
M_c=L_c+R_c.
}
\]

\begin{proposition}
\[
\det M_c>0
\iff
\text{every connected group of color }c\text{ has at least one liberty}.
\]
\end{proposition}

\begin{proof}
For any real vector \(x\) indexed by stones of color \(c\),
\[
x^TM_cx
=
\sum_{\langle ij\rangle\in G_c}(x_i-x_j)^2
+
\sum_i r_i x_i^2.
\]

Every term is nonnegative. Therefore \(x^TM_cx=0\) requires \(x_i=x_j\) along every same-color edge, so \(x\) is constant on every connected component.

The second term additionally requires \(x_i=0\) wherever \(r_i>0\). Thus a nonzero zero-mode exists precisely when some connected component contains no vertex with \(r_i>0\).

But \(r_i>0\) means that stone touches an empty site. Hence \(M_c\) is nonsingular exactly when every group possesses a liberty.
\end{proof}

\section{Why the determinant is more than an algebraic trick}

The same determinant has a combinatorial meaning through the matrix-forest theorem.

Add a ghost vertex \(\star\). Connect every stone \(i\) to \(\star\) with edge weight \(r_i\).

A spanning tree of the enlarged graph becomes, after deleting \(\star\), a forest in which every tree component is attached to one permitted root.

In our interpretation,
\[
\boxed{
\text{every stone follows a forest path toward a liberty.}
}
\]

For a single isolated stone,
\[
L=[0],
\qquad
R=[r],
\]
so
\[
\det(L+R)=r.
\]

For two connected stones,
\[
L=
\begin{pmatrix}
1&-1\\
-1&1
\end{pmatrix},
\qquad
R=
\begin{pmatrix}
r_1&0\\
0&r_2
\end{pmatrix},
\]
and therefore
\[
\det(L+R)=r_1+r_2+r_1r_2.
\]

Again, the determinant vanishes only when both stones lack liberties.

\section{An exact statistical-mechanical partition function}

Now define
\[
D_c[s]=\det M_c[s].
\]

Because \(M_c\) is positive semidefinite,
\[
D_c[s]\geq0.
\]

Consider
\[
\boxed{
Z_\alpha
=
\sum_{\{s_i\}}
e^{-\beta H_0[s]}
\prod_{c=\pm1}D_c[s]^\alpha.
}
\]

Call this provisionally the Go--spanning-forest model.

As
\[
\alpha\rightarrow0^+,
\]
for a legal configuration \(D_c>0\), hence
\[
D_c^\alpha\rightarrow1,
\]
while for an illegal configuration \(D_c=0\), hence
\[
D_c^\alpha=0.
\]

Thus
\[
\boxed{
\lim_{\alpha\to0^+}
\prod_cD_c^\alpha
=
\mathbf 1_{\rm all\ groups\ have\ liberties}.
}
\]

Consequently,
\[
\boxed{
Z_{0^+}
=
\sum_{\text{liberty-legal }s}
e^{-\beta H_0[s]}.
}
\]

At \(\alpha=1\), legal boards receive a weight proportional to the number of rooted-forest certificates connecting their stones to liberties.

For general \(\alpha\),
\[
F_{\rm forest}
=
-\frac{\alpha}{\beta}
\sum_c\log\det M_c
\]
acts as a liberty entropy.

\section{The connection to spanning-forest field theory}

For Grassmann variables \(\psi_i,\bar\psi_i\),
\[
\det M
=
\int
D\bar\psi\,D\psi\,
e^{-\bar\psi M\psi}.
\]

Therefore the \(\alpha=1\) model can be written schematically as
\[
Z_1
=
\sum_s
\int
D\bar\psi_+D\psi_+
D\bar\psi_-D\psi_-
\,
e^{-S[s,\psi_\pm]},
\]
with
\[
S
=
\beta H_0[s]
+
\bar\psi_+M_+[s]\psi_+
+
\bar\psi_-M_-[s]\psi_-.
\]

Thus the spin-1 lattice gas becomes coupled to fermionic auxiliary fields.

In a large stone domain \(\Omega_c\), one expects a coarse-grained operator resembling
\[
M_c
\rightsquigarrow
-\nabla^2_{\Omega_c}
\]
with absorbing or mixed boundary conditions on portions of the boundary that touch liberties.

A continuum interfacial free energy may therefore contain
\[
F[\Omega_c]
=
\sigma|\partial\Omega_c|
-\Delta f|\Omega_c|
-\alpha T\log\det(-\nabla^2_{\Omega_c})
+\cdots.
\]

\section{The empty-space complex}

Let \(E_s\) denote the graph consisting of empty lattice sites and nearest-neighbor edges between them.

Its zeroth cohomology is
\[
H^0(E_s;\Bbbk).
\]

For a finite graph,
\[
\dim H^0(E_s;\Bbbk)
=
\beta_0(E_s)
=
\#\{\text{connected empty regions}\}.
\]

If
\[
E_s=D_1\sqcup\cdots\sqcup D_r,
\]
then one may choose the canonical indicator cocycles
\[
\chi_{D_a}(i)
=
\begin{cases}
1,&i\in D_a,\\
0,&i\notin D_a.
\end{cases}
\]

Each
\[
[\chi_{D_a}]\in H^0(E_s)
\]
represents one empty chamber.

\section{The physical boundary and the barrier complex}

A Go territory may lie against the edge or in a corner. The physical boundary of the board acts as part of the enclosure.

Represent the board as a topological disk
\[
X\simeq D^2.
\]

Thicken stones slightly into two-dimensional cells and let
\[
K_s
\]
denote their union.

Now define the barrier complex
\[
\boxed{
B_s=K_s\cup\partial X.
}
\]

The connected components of
\[
X\setminus K_s
\]
are the empty chambers.

Because the physical boundary \(\partial X\) belongs to \(B_s\), a stone arc connecting two portions of the board edge can generate a closed barrier together with the intervening piece of \(\partial X\).

\section{Alexander duality: why loops create chambers}

For a compact, sufficiently regular subset \(B\subset S^2\), Alexander duality gives
\[
\widetilde H_i(S^2-B;\Bbbk)
\cong
\widetilde H^{1-i}(B;\Bbbk).
\]

For \(i=0\),
\[
\boxed{
\widetilde H_0(S^2-B)
\cong
\widetilde H^1(B).
}
\]

Thus additional connected components of the complement correspond to first cohomology classes of the barrier.

Applied to
\[
B_s=K_s\cup\partial X,
\]
the statement says approximately
\[
\boxed{
\text{independent empty chambers}
\longleftrightarrow
\text{independent }H^1\text{ classes of the barrier}.
}
\]

A convenient enclosure count is
\[
\mathcal E(s)=\beta_1(B_s)-1,
\]
where the subtraction removes the trivial cycle supplied by the board edge.

\section{Ownership is topology plus a color boundary condition}

For an empty component \(D\), define the set of colors of neighboring stones,
\[
C(D)
=
\{s_i\in\{-1,+1\}:
i\text{ is adjacent to some }j\in D\}.
\]

Then define
\[
o(D)
=
\begin{cases}
+1,&C(D)=\{+1\},\\
-1,&C(D)=\{-1\},\\
0,&C(D)=\{-1,+1\},\\
0,&C(D)=\varnothing.
\end{cases}
\]

Let
\[
\Omega_E=\sum_{i:s_i=0}[i]
\]
be the empty-site 0-chain. Then
\[
\langle\chi_D,\Omega_E\rangle
=
|D|.
\]

A natural topological territory functional is
\[
\boxed{
T[s]
=
\sum_{D\in\pi_0(E_s)}
o(D)
\langle\chi_D,\Omega_E\rangle.
}
\]

Topology determines which regions exist, measure determines how large they are, and color boundary data determines who owns them.

\section{Eyes as complement topology}

Consider one connected stone group \(G\). Look at empty components whose stone boundary belongs entirely to \(G\).

Define schematically
\[
N_{\rm eye}^{\rm top}(G)
=
\#\{
D\subset E:
\partial_{\rm stones}D\subseteq G
\}.
\]

A group with two disjoint internal empty regions has
\[
N_{\rm eye}^{\rm top}\geq2.
\]

But a false eye may be a genuine geometric hole and still fail tactically. Therefore
\[
\boxed{
\text{topological eye}
\neq
\text{dynamically protected eye}.
}
\]

\section{The minimal physics benchmark: the Ising contour gas}

Take ordinary Ising spins
\[
\tau_i=\pm1
\]
with
\[
H_{\rm Ising}
=
-J\sum_{\langle ij\rangle}\tau_i\tau_j.
\]

Fix \(+\) boundary conditions. For every bond joining opposite spins, draw an edge on the dual lattice.

These dual edges form closed Peierls contours. The low-temperature partition function can be written
\[
\boxed{
Z_+
=
e^{\beta JN_b}
\sum_{\Gamma}
e^{-2\beta J|\Gamma|},
}
\]
where \(\Gamma\) runs over compatible closed contour configurations and \(|\Gamma|\) is total contour length.

This gives an exact demonstration of
\[
\boxed{
\text{nearest-neighbor local Hamiltonian}
\longrightarrow
\text{extended closed loops}.
}
\]

\section{Perimeter versus area: territory as droplet physics}

Add a magnetic field,
\[
H
=
-J\sum_{\langle ij\rangle}\tau_i\tau_j
-h\sum_i\tau_i.
\]

A coarse-grained droplet free energy has the form
\[
\boxed{
\Delta F(D)
=
\sigma P(D)-\Delta f A(D).
}
\]

For a circular droplet,
\[
P=2\pi R,
\qquad
A=\pi R^2,
\]
so
\[
\Delta F(R)
=
2\pi\sigma R-\pi\Delta fR^2.
\]

The critical radius is
\[
R_c=\frac{\sigma}{\Delta f}.
\]

Thus
\[
\boxed{
\text{boundary cost}\sim R,
\qquad
\text{interior reward}\sim R^2.
}
\]

\section{Relation to \(O(n)\) loop gases}

A standard loop model has partition function
\[
Z_{\rm loop}
=
\sum_{\Gamma}
x^{|\Gamma|}
n^{N_{\rm loop}(\Gamma)}.
\]

Our barrier model suggests introducing
\[
y^{\mathcal E(s)}
=
y^{\beta_1(B_s)-1}.
\]

An analytically useful deformation is
\[
Z_{\alpha,y}
=
\sum_s
e^{-\beta H_0[s]}
\prod_cD_c[s]^\alpha
y^{\beta_1(B_s)-1}.
\]

The important experiment is to start with \(y=1\), so enclosure has no explicit energetic reward. If the system nevertheless develops macroscopic \(\beta_1\) structure, it is genuinely emergent.

\section{Relation to the six-vertex model}

The six-vertex or ice model places arrows on lattice edges and imposes the purely local rule
\[
2\text{ arrows in},
\qquad
2\text{ arrows out}
\]
at every vertex.

A local divergence-free constraint generates conserved fluxes, loop structures, and height fields.

Conceptually,
\[
\nabla\cdot E=0
\]
locally forces extended global structure.

This is close in spirit to the present construction, except that our fundamental constraint is the existence of a rooted path from every stone to a liberty.

\section{Cohomology versus topological order}

Suppose a binary site field \(u\in C^0(X;\mathbb Z_2)\) describes two phases.

Its domain-wall field is roughly
\[
a=du\in C^1(X;\mathbb Z_2).
\]

Automatically,
\[
da=d^2u=0.
\]

So the wall field is a cocycle, but it is exact:
\[
[a]=0\in H^1(X;\mathbb Z_2).
\]

Thus merely obtaining closed loops does not imply a nontrivial cohomology sector.

If we want genuine gauge-like topological sectors, the edge variable must instead be promoted to an independent field
\[
a\in C^1(X;\mathbb Z_2)
\]
subject only to
\[
da=0,
\]
without requiring \(a=du\).

On
\[
X=T^2,
\]
we have
\[
H^1(T^2;\mathbb Z_2)
=
\mathbb Z_2\oplus\mathbb Z_2.
\]

The two bits record winding around the two noncontractible cycles.

Hence
\[
\boxed{
\begin{array}{c}
\text{closed territorial loops}\\
\neq\\
\text{nontrivial cohomology sector}\\
\neq\\
\text{topological order}.
\end{array}
}
\]

\section{Mean-field solution of the unconstrained spin sector}

Define
\[
m=\langle s_i\rangle,
\qquad
q=\langle s_i^2\rangle.
\]

For coordination number \(z=4\), mean-field decoupling gives the single-site effective energy
\[
\epsilon(s)
=
-(zJm+h)s
+
(\Delta-zKq)s^2.
\]

Define
\[
u=\beta(zJm+h)
\]
and
\[
a=e^{-\beta(\Delta-zKq)}.
\]

The one-site partition function is
\[
Z_1
=
1+2a\cosh u.
\]

Therefore
\[
\boxed{
m
=
\frac{2a\sinh u}
{1+2a\cosh u}
}
\]
and
\[
\boxed{
q
=
\frac{2a\cosh u}
{1+2a\cosh u}.
}
\]

The mean-field free energy is
\[
\boxed{
f(m,q)
=
\frac{zJ}{2}m^2
+
\frac{zK}{2}q^2
-
\frac1\beta
\log\!\left[
1+
2e^{-\beta(\Delta-zKq)}
\cosh\!\big(\beta(zJm+h)\big)
\right].
}
\]

At \(h=0\), linearizing near \(m=0\) gives
\[
m
\simeq
\beta zJ q\,m,
\]
so the mean-field instability occurs when
\[
\boxed{
1=\beta zJq.
}
\]

\section{What the forest constraint changes}

Write the full effective free energy as
\[
H_{\rm eff}
=
H_0
-
\frac{\alpha}{\beta}
\sum_c\log\det M_c.
\]

In eigenvalues,
\[
\log\det M_c
=
\sum_a\log\lambda_a^{(c)}.
\]

This suggests defining a liberty gap
\[
\boxed{
\Delta_{\rm lib}(C)
=
\lambda_{\min}(L_C+R_C).
}
\]

A well-rooted compact group should have a relatively large liberty gap. A group connected to its only liberty through a long narrow corridor will tend to have a much smaller gap. A dead group has
\[
\Delta_{\rm lib}=0.
\]

\section{The proposed phase diagram}

The natural control parameters are
\[
\beta J,\qquad
\beta K,\qquad
\beta\Delta,\qquad
\alpha.
\]

Measure
\[
m,\quad
q,\quad
\beta_0(E_s),\quad
\beta_1(B_s),\quad
A_{\max}^{\rm owned},\quad
\Delta_{\rm lib}.
\]

Use as a primary territorial order parameter
\[
\boxed{
\rho_{\rm terr}
=
\frac{1}{L^2}
\mathbb E
\left[
\max_{D:\,o(D)\neq0}|D|
\right].
}
\]

A genuine macroscopic territorial phase would have
\[
\lim_{L\rightarrow\infty}\rho_{\rm terr}>0.
\]

Define a topological susceptibility
\[
\chi_{\rm top}
=
\frac{1}{L^2}
\left(
\langle\mathcal E^2\rangle
-
\langle\mathcal E\rangle^2
\right).
\]

\section{The key experiment: does topology arise without rewarding it?}

The trivial model is
\[
H=H_{\rm local}-gA_{\rm territory}.
\]

If \(g>0\), the system is explicitly rewarded for territory.

The interesting model is
\[
\boxed{
H
=
H_0
-
\frac{\alpha}{\beta}
\sum_c\log\det M_c,
}
\]
with no explicit area or topology term.

If macroscopic enclosure appears spontaneously, then
\[
\boxed{
\text{local energetic interactions}
+
\text{local-certifiable survival constraint}
\rightarrow
\text{emergent topology of the complement}.
}
\]

\section{Go is fundamentally dynamical}

A Go move is a map
\[
s_t\rightarrow s_{t+1}.
\]

Suppose color \(c_t\) plays at empty site \(x\).

First perform placement,
\[
s_x\leftarrow c_t.
\]

Then inspect adjacent opponent groups. Any opponent group \(C\) satisfying
\[
L(C)=0
\]
is removed:
\[
s_i\leftarrow0,
\qquad i\in C.
\]

After opponent captures are removed, inspect the newly placed player's own group. If it has no liberty, reject the move under ordinary no-suicide rules.

Finally impose ko or superko.

Schematically,
\[
\boxed{
\mathcal T_{c,x}
=
\mathcal P_{\rm ko}
\,
\mathcal P_{\rm suicide}
\,
\mathcal C_{\rm capture}
\,
\mathcal P_{c,x}.
}
\]

\section{A driven physical model inspired by Go}

Define a stochastic lattice process.

Black and White alternately attempt depositions. At an empty site \(x\), color \(c\) attempts
\[
0\rightarrow c
\]
with rate
\[
r_c(x|s)
=
r_0
\min\left[
1,
e^{-\beta\Delta H_0}
\right].
\]

After deposition, zero-liberty enemy clusters are removed. A self-killing move is rejected.

Its generator can be written schematically as
\[
(\mathcal L f)(s)
=
\sum_{c,x}
r_c(x|s)
\left[
f(\mathcal T_{c,x}s)-f(s)
\right].
\]

Because capture performs irreversible many-site changes, detailed balance generally fails.

\section{Capture avalanches}

Let
\[
S_{\rm cap}
\]
be the number of stones removed in one capture event.

Measure
\[
P(S_{\rm cap}).
\]

Near a dynamical critical regime, one might find
\[
P(S_{\rm cap})
\sim
S_{\rm cap}^{-\tau}
f(S_{\rm cap}/S_c).
\]

Also measure
\[
\Delta\beta_1
=
\beta_1(B_{t+1})
-
\beta_1(B_t),
\]
and the joint distribution
\[
P(S_{\rm cap},\Delta\beta_1).
\]

\section{Time-dependent topology and zigzag persistence}

Ordinary persistent homology assumes
\[
X_0\subseteq X_1\subseteq X_2\subseteq\cdots.
\]

Go does not.

One can instead form
\[
K_t
\hookrightarrow
K_t\cup K_{t+1}
\hookleftarrow
K_{t+1},
\]
giving
\[
K_0
\rightarrow
K_0\cup K_1
\leftarrow
K_1
\rightarrow
K_1\cup K_2
\leftarrow
K_2
\rightarrow\cdots.
\]

Applying \(H^1\) gives a zigzag persistence module. Alternatively apply \(H^0\) to empty-space graphs.

The resulting barcode records the lifetimes of enclosed regions.

Define, for example,
\[
P_{\rm top}(s_{0:T})
=
\sum_a
\ell_a^p,
\]
where \(\ell_a\) is the lifetime of the \(a\)-th enclosure class.

\section{A definition of settled territory}

Call a chamber class \(\tau\)-persistent if its zigzag interval survives for at least \(\tau\) future moves.

Define
\[
A_{\rm stable}(t,\tau)
=
\sum_{\alpha:\ell_\alpha\geq\tau}
A(\alpha).
\]

Then
\[
\frac{A_{\rm stable}(t,\tau)}{L^2}
\]
is a quantitative measure of how much of the board has become topologically settled.

\section{The strategic model}

Actual Go introduces adversarial optimization.

Let a terminal score be
\[
S[s]
=
A_B[s]-A_W[s]
+
N_B[s]-N_W[s]
-\text{komi},
\]
with the precise form adjusted for scoring convention.

A perfect zero-sum value function satisfies
\[
V_c(s)
=
\max_{x\in\mathcal L_c(s)}
\left[
-V_{-c}(\mathcal T_{c,x}s)
\right],
\]
with
\[
V(s_{\rm terminal})=S[s_{\rm terminal}].
\]

Compare three ensembles:
\[
\text{equilibrium},
\qquad
\text{driven stochastic},
\qquad
\text{adversarial strategic}.
\]

If only strategic dynamics develops macroscopic enclosure, then
\[
\boxed{
\text{topological organization is strategy-induced rather than equilibrium-induced}.
}
\]

\section{Relation to random-cluster and forest physics}

The Fortuin--Kasteleyn random-cluster model unifies percolation, Ising, and Potts physics.

The arboreal gas appears in an appropriate
\[
q\rightarrow0
\]
limit, where graph configurations become forests.

The present construction naturally couples a spin-1 occupancy/color sector to a rooted-forest sector:
\[
\boxed{
\text{BEG spin-1 model}
\quad\text{coupled to}\quad
\text{rooted forest / }q\rightarrow0\text{ sector}.
}
\]

\section{Model correspondence table}

\begin{center}
\begin{tabular}{p{0.28\textwidth}p{0.31\textwidth}p{0.31\textwidth}}
\toprule
Our ingredient & Established analogue & Shared structure\\
\midrule
\(s_i=-1,0,+1\) & Blume--Emery--Griffiths / Blume--Capel & two species plus vacancy\\
Friendly/enemy interfaces & Ising/Potts & local domain energetics\\
Closed boundaries & Peierls contours & local interactions become loops\\
Loop fugacity & \(O(n)\) model & weight per extended loop\\
Local conservation constraint & six-vertex / ice & local rule produces global flux structure\\
Liberty condition & rooted spanning forests & every component connects to a permitted root\\
\(\det(L+R)\) & matrix-forest theorem & determinant counts rooted forests\\
Forest continuum sector & \(q\to0\) Potts / arboreal gas & fermionic and supersymmetric formulations\\
Capture dynamics & bootstrap/KCM-like culling & local constraints generate complex dynamics\\
Territory chambers & complement \(H^0\) & connected regions are topological classes\\
Enclosing barriers & \(H^1\) via Alexander duality & loops dual to complement components\\
Time-dependent territory & zigzag persistence & topology under additions and deletions\\
Torus extension & lattice gauge theory & nontrivial cohomology/winding sectors\\
\bottomrule
\end{tabular}
\end{center}

\section{How to solve the model numerically}

For each proposed configuration \(s\), construct the connected same-color graphs \(G_+\) and \(G_-\).

Union-find gives their components essentially in linear time.

The empty graph similarly yields the chamber decomposition.

The root matrices are obtained from
\[
r_i=\lambda\ell_i.
\]

For a \(19\times19\) board, the matrices have at most 361 vertices, so direct sparse factorizations are feasible.

Compute
\[
\log D_c
=
\log\det(L_c+R_c)
\]
using sparse Cholesky decomposition whenever \(M_c\) is positive definite.

The spin sector can be sampled with local Metropolis updates, cluster-inspired moves, or joint spin/forest updates.

\section{Finite-size scaling}

For
\[
L=16,24,32,48,64,\ldots,
\]
measure
\[
m,\quad
q,\quad
\rho_{\rm terr},\quad
\mathcal E,\quad
\chi_{\rm top},\quad
\Delta_{\rm lib},
\]
and loop/chamber size distributions.

Near a continuous transition one expects
\[
\rho_{\rm terr}(L,t)
=
L^{-\beta_T/\nu}
F(tL^{1/\nu}),
\]
where
\[
t=\frac{g-g_c}{g_c}.
\]

Similarly,
\[
\chi_{\rm top}
\sim
L^{\gamma_T/\nu}.
\]

At a candidate conformal critical point, transfer-matrix calculations on long strips could extract an effective central charge from
\[
f_L
=
f_\infty
-\frac{\pi c_{\rm eff}}{6L^2}
+\cdots.
\]

\section{A clean analytically tractable hierarchy}

The zeroth model is ordinary Ising.

The next model is the BEG lattice gas without forest constraints.

Then switch on a weak forest exponent,
\[
0<\alpha\ll1.
\]

Formally,
\[
D^\alpha
=
e^{\alpha\log D}
=
1+\alpha\log D+O(\alpha^2).
\]

Therefore
\[
H_{\rm eff}
=
H_0
-\frac{\alpha}{\beta}\log D_+
-\frac{\alpha}{\beta}\log D_-
+O(\alpha^2).
\]

At
\[
\alpha=1,
\]
the auxiliary fermion/forest representation becomes especially natural.

\section{A possible continuum effective theory}

Introduce coarse-grained fields
\[
\rho(x)
\]
for occupation density and
\[
\phi(x)
\]
for color polarization.

A standard Landau functional is
\[
S_{\rm BEG}
=
\int d^2x\,
\left[
\frac{\kappa_\phi}{2}(\nabla\phi)^2
+
\frac{\kappa_\rho}{2}(\nabla\rho)^2
+
V(\phi,\rho)
\right].
\]

The forest fields introduce Grassmann variables
\[
\psi_\pm,\bar\psi_\pm.
\]

Schematically,
\[
S_{\rm forest}
=
\int d^2x\,
\bar\psi_+
\mathcal M_+[\phi,\rho]
\psi_+
+
\bar\psi_-
\mathcal M_-[\phi,\rho]
\psi_-.
\]

A minimal operator might have the form
\[
\mathcal M_\pm
=
-\nabla\cdot
D_\pm(\phi,\rho)
\nabla
+
\mu_\pm(\phi,\rho,\nabla\rho).
\]

After integrating out the fermions,
\[
S_{\rm eff}
=
S_{\rm BEG}
-
\alpha\log\det\mathcal M_+
-
\alpha\log\det\mathcal M_-.
\]

\section{What could actually be proved}

Several pieces are theorem-level statements:

\begin{enumerate}
\item The liberty-determinant theorem:
\[
\det(L_c+R_c)>0
\]
if and only if each color-\(c\) component possesses a liberty.

\item The rooted-forest interpretation of that determinant via the matrix-forest theorem.

\item The planar duality statement relating complement components to \(H^1\) of the barrier through Alexander duality.

\item The exact Peierls contour representation of the Ising benchmark.
\end{enumerate}

The central open statements concern the coupled model.

\begin{conjecture}[Enclosure Transition]
There exists a region of parameter space of the Go--spanning-forest model in which a thermodynamic transition separates a phase with only microscopic owned chambers from a phase in which
\[
\lim_{L\to\infty}
\frac{\mathbb E[A_{\max}^{\rm owned}]}{L^2}>0.
\]
\end{conjecture}

\begin{conjecture}[Driven Enclosure]
Alternating deposition-and-capture dynamics generates a nonequilibrium regime with macroscopic persistent complement classes despite the absence of an explicit topological energy term.
\end{conjecture}

\begin{conjecture}[Strategy-Induced Topology]
Adversarial optimization produces a substantially different distribution of persistent \(H^1\) classes from both equilibrium and unbiased driven dynamics.
\end{conjecture}

\section{Why the forest connection may be the deepest part}

There is a conceptual triangle:
\[
\boxed{
\begin{array}{ccc}
\text{connectivity} &\longleftrightarrow& \text{graph Laplacian}\\
\updownarrow && \updownarrow\\
\text{survival to liberty} &\longleftrightarrow& \det(L+R)\\
\updownarrow && \updownarrow\\
\text{rooted forest} &\longleftrightarrow& \text{fermionic field theory}.
\end{array}
}
\]

The Laplacian controls diffusion, random walks, electrical networks, spectral geometry, Gaussian fields, spanning trees, and many continuum field theories.

Liberties become absorbing roots.

A Go group becomes a domain in which a random walker must eventually be able to reach an absorbing boundary.

A dead group is precisely a component with no absorbing root and hence a Laplacian zero mode.

Thus
\[
\boxed{
\text{A living group is a domain with a sink.}
}
\]

Capture corresponds to the spectral singularity
\[
\lambda_{\min}\to0,
\qquad
\det M\to0.
\]

\section{A probabilistic interpretation}

Because \(M=L+R\) is a killed graph Laplacian, consider a random walk moving among stones of one group.

Whenever it reaches a root site, it can escape through a liberty.

The Green function
\[
G=(L+R)^{-1}
\]
measures propagation before killing.

As the group approaches death,
\[
\lambda_{\min}\to0,
\]
so \(G\) develops a large mode.

This suggests vulnerability measures such as
\[
\boxed{
\mathcal V(C)
=
\operatorname{Tr}(L_C+R_C)^{-1}
}
\]
or
\[
\mathcal V_\infty(C)
=
\frac1{\lambda_{\min}}.
\]

\section{Where sheaf cohomology might enter}

Ordinary cellular cohomology is sufficient for enclosure.

A sheaf formulation becomes meaningful only if each local patch \(U\) carries a space
\[
\mathcal F(U)
\]
of locally consistent tactical claims, with restriction maps
\[
\rho_{UV}:\mathcal F(U)\to\mathcal F(V).
\]

A global consistent interpretation is a section
\[
s\in H^0(X;\mathcal F),
\]
while a nontrivial
\[
H^1(X;\mathcal F)
\]
would represent an obstruction to gluing locally consistent tactical assessments into a globally consistent one.

Without a precise local data structure and restriction map, sheaf language should not be introduced.

\section{A toroidal extension}

Put the model on
\[
T^2.
\]

Introduce an independent closed edge field
\[
a\in Z^1(T^2;\mathbb Z_2).
\]

Its cohomology class is
\[
[a]=(w_x,w_y)
\in
H^1(T^2;\mathbb Z_2)
\simeq
\mathbb Z_2^2.
\]

Define Wilson-type observables
\[
W_x=(-1)^{w_x},
\qquad
W_y=(-1)^{w_y}.
\]

Then study sector-dependent free energies
\[
F_{w_x,w_y}(L).
\]

This is the appropriate path toward genuine gauge/topological physics.

\section{A realistic first-paper program}

Define
\[
\boxed{
Z_\alpha(J,K,\Delta)
=
\sum_s
e^{-\beta H_{\rm BEG}[s]}
\prod_{c=\pm}
\det\!\left(
L_c[s]+\lambda R_c[s]
\right)^\alpha.
}
\]

Then:

\begin{enumerate}
\item Prove the determinant/liberty theorem.
\item Derive the rooted-forest representation.
\item Develop mean-field theory for the BEG sector.
\item Simulate the full model.
\item Measure conventional and topological observables.
\item Determine whether the forest constraint shifts known transition lines or creates qualitatively new phases.
\item Search for a regime with
\[
m=0
\]
but
\[
\rho_{\rm terr}>0.
\]
\item If a continuous transition is found, perform finite-size scaling and compare against Ising, percolation, Potts, and loop-model universality classes.
\end{enumerate}

\section{The second paper}

Replace equilibrium updates with Go-like deposition/capture dynamics:
\[
s_{t+1}
=
\mathcal T_{c_t,x_t}s_t.
\]

Study
\[
P(S_{\rm cap}),
\quad
P(\Delta\beta_1),
\quad
\text{activity},
\quad
\text{relaxation time},
\quad
\text{persistent territory}.
\]

Compare this driven ensemble to an equilibrium ensemble having similar one-point densities.

\section{The third paper}

Use actual game records.

For every move \(t\), construct
\[
E_t,\qquad
B_t,\qquad
L_\pm(t)+R_\pm(t).
\]

Measure
\[
\beta_0(E_t),
\qquad
\beta_1(B_t),
\qquad
\lambda_{\min}^{\pm}(t),
\]
and zigzag persistence intervals.

Compare professional games, strong AI, weak AI, and random legal play.

\section{Conceptual endpoint}

The project develops through the chain
\[
\begin{aligned}
\text{site variables}
&\longrightarrow
\text{same-color graph}\\
&\longrightarrow
\text{graph Laplacian}\\
&\longrightarrow
\text{rooted forest}\\
&\longrightarrow
\text{liberty constraint}\\
&\longrightarrow
\text{surviving extended barriers}\\
&\longrightarrow
H^1(B)\\
&\overset{\rm Alexander}{\longleftrightarrow}
H^0(X-B)\\
&\longrightarrow
\text{enclosed chambers}\\
&\longrightarrow
\text{territorial area}.
\end{aligned}
\]

The dynamics adds
\[
\text{capture}
\longrightarrow
\text{birth/death of topological classes}
\longrightarrow
\text{zigzag persistence}.
\]

The field theory adds
\[
\det(L+R)
\longrightarrow
\text{Grassmann variables}
\longrightarrow
\text{forest/supersymmetric sector}.
\]

The gauge extension adds
\[
da=0,
\qquad
[a]\in H^1(T^2;\mathbb Z_2).
\]

Thus a better summary than ``Go is an Ising model'' is
\[
\boxed{
\text{Go-like physics}
=
\text{spin-1 lattice gas}
+
\text{rooted-forest constraints}
+
\text{loop/complement topology}
+
\text{nonequilibrium adversarial dynamics}.
}
\]

The central mathematical object is
\[
\boxed{
Z_\alpha
=
\sum_{\{s_i=-1,0,+1\}}
e^{-\beta H_{\rm BEG}[s]}
\prod_{c=\pm}
\det(L_c[s]+R_c[s])^\alpha.
}
\]

The clean research question is therefore
\[
\boxed{
\textbf{
Can a locally interacting, locally certifiable lattice system undergo a transition in the topology of its complement?
}
}
\]

Go is the motivating system that tells us which local constraint to study.

\end{document}
